\section{Phản ứng trung hòa}

\subsection{Bản chất của phản ứng}

Phản ứng giữa acid và base tạo muối và nước được gọi là \textbf{phản ứng trung hòa}. Với acid mạnh và base mạnh, có thể hình dung phần cốt lõi của phản ứng là:
\[
\ce{H+ + OH- -> H2O}.
\]

Ví dụ:
\[
\ce{HCl + NaOH -> NaCl + H2O},
\]
\[
\ce{H2SO4 + 2KOH -> K2SO4 + 2H2O}.
\]

\begin{ghinho}{Các bước giải bài trung hòa}
\begin{enumerate}
\item Tính số mol các chất từ dữ kiện đề bài.
\item Viết và cân bằng phương trình hóa học.
\item Tính các giá trị trung gian cần thiết: chất hết, chất dư, số mol sản phẩm,\ldots
\item Tính đại lượng đề hỏi và kết luận bằng câu đầy đủ.
\end{enumerate}
\end{ghinho}

\begin{vidu}
Trộn \SI{100}{\milli\liter} dung dịch \ce{HCl} \SI{0.50}{\mole\per\liter} với \SI{150}{\milli\liter} dung dịch \ce{NaOH} \SI{0.20}{\mole\per\liter}. Xác định chất còn dư sau phản ứng.
\end{vidu}
\begin{loigiai}
\textbf{Bước 1. Tính số mol.}
\[
n_{\ce{HCl}}=0.100\times0.50=\SI{0.050}{\mole},
\]
\[
n_{\ce{NaOH}}=0.150\times0.20=\SI{0.030}{\mole}.
\]

\textbf{Bước 2. Viết phương trình hóa học và đặt số mol vào phương trình.}
\begin{center}\small
\setlength{\tabcolsep}{3pt}
\begin{tabular}{c c c c c c c}
\ce{HCl} & + & \ce{NaOH} & \ce{->} & \ce{NaCl} & + & \ce{H2O}\\
$0{,}050$ & & $0{,}030$ & & & & \\
$0{,}030$ & & $0{,}030$ & & $0{,}030$ & & $0{,}030$
\end{tabular}
\end{center}

\textbf{Bước 3. Tính các giá trị trung gian.}

Theo phương trình, \ce{NaOH} phản ứng hết; số mol \ce{HCl} đã phản ứng là \SI{0.030}{\mole}.
\[
n_{\ce{HCl}\,\text{dư}}=0.050-0.030=\SI{0.020}{\mole}.
\]

\textbf{Bước 4. Tính đại lượng đề hỏi và kết luận.}

\textbf{Vậy,} sau phản ứng còn dư \SI{0.020}{\mole} \ce{HCl}.
\end{loigiai}

\subsection{Trung hòa trong tình huống thực tế}

Nhiều bài toán không nói trực tiếp ``acid phản ứng với base'' mà đặt trong bối cảnh: thuốc dạ dày, xử lí đất chua, xử lí nước thải, làm sạch acid dư. Cách giải vẫn không đổi: xác định chất nào đóng vai trò acid, chất nào đóng vai trò base rồi viết phản ứng.

\begin{vidu}
Một loại thuốc dạ dày chứa \ce{Mg(OH)2}. Một liều thuốc có \SI{0.58}{\gram} \ce{Mg(OH)2}. Tính số mol \ce{HCl} tối đa có thể trung hòa. Cho $M_{\ce{Mg(OH)2}}=58$ g/mol.
\end{vidu}
\begin{loigiai}
\textbf{Bước 1. Tính số mol.}
\[
n_{\ce{Mg(OH)2}}=\frac{0.58}{58}=\SI{0.010}{\mole}.
\]

\textbf{Bước 2. Viết phương trình hóa học và đặt số mol vào phương trình.}
\begin{center}\small
\setlength{\tabcolsep}{3pt}
\begin{tabular}{c c c c c c c}
\ce{Mg(OH)2} & + & \ce{2HCl} & \ce{->} & \ce{MgCl2} & + & \ce{2H2O}\\
$0{,}010$ & & $0{,}020$ & & $0{,}010$ & & $0{,}020$
\end{tabular}
\end{center}

\textbf{Bước 3. Tính các giá trị trung gian.}
\[
n_{\ce{HCl}}=2\times0.010=\SI{0.020}{\mole}.
\]

\textbf{Bước 4. Tính đại lượng đề hỏi và kết luận.}

\textbf{Vậy,} một liều thuốc có thể trung hòa tối đa \SI{0.020}{\mole} \ce{HCl}.
\end{loigiai}

\begin{tuluyen}
Trộn \SI{200}{\milli\liter} dung dịch \ce{H2SO4} \SI{0.10}{\mole\per\liter} với \SI{150}{\milli\liter} dung dịch \ce{NaOH} \SI{0.30}{\mole\per\liter}. Xác định chất dư và số mol chất dư.
\dapso{\ce{NaOH} dư \SI{0.005}{\mole}.}
\begin{loigiaiBT}
\textbf{Bước 1. Tính số mol.}
\[
n_{\ce{H2SO4}}=0.200\times0.10=\SI{0.020}{\mole},
\]
\[
n_{\ce{NaOH}}=0.150\times0.30=\SI{0.045}{\mole}.
\]

\textbf{Bước 2. Viết phương trình hóa học và đặt số mol vào phương trình.}
\begin{center}\small
\setlength{\tabcolsep}{3pt}
\begin{tabular}{c c c c c c c}
\ce{H2SO4} & + & \ce{2NaOH} & \ce{->} & \ce{Na2SO4} & + & \ce{2H2O}\\
$0{,}020$ & & $0{,}045$ & & & & \\
$0{,}020$ & & $0{,}040$ & & $0{,}020$ & & $0{,}040$
\end{tabular}
\end{center}

\textbf{Bước 3. Tính các giá trị trung gian.}
\[
n_{\ce{NaOH}\,\text{dư}}=0.045-0.040=\SI{0.005}{\mole}.
\]

\textbf{Bước 4. Tính đại lượng đề hỏi và kết luận.}

\textbf{Vậy,} sau phản ứng còn dư \SI{0.005}{\mole} \ce{NaOH}.
\end{loigiaiBT}
\end{tuluyen}
